\section{Background}
\label{sec:background}

\subsection{Words and Trees}\label{sec:words}
Fix an arity $n\in\mathbb N$, the child alphabet $\Dir=\{0,\ldots,n-1\}$,
and a variance map
$c:\Dir\to\{0,1\}$. Here $0$ means covariant and $1$ contravariant.
The signature contains exactly $\Bot$, $\Top$, and one constructor $f$ of arity $n$.
A \emph{cut} of a word $w$ of length $m$ is an integer $0\le j\le m$;
cut $j$ separates its first $j$ letters from the remaining suffix.
Write $w[i]$ for its letter at position $i$, with indices starting at $0$
and $0\le i<m$. Words are finite; $\eps$ is the empty word, concatenation is juxtaposition,
$u\pre w$ means that $u$ is a prefix of $w$, and $w[:j]$, $w[j:h]$, $w[j:]$
are its prefix, factor, and suffix at the indicated cuts. For example, over
$\{0,1\}$, $010[:2]=01$ and $010[1:]=10$. A positive integer $\varpi$ is a
\emph{period} of $w$ when $w[i]=w[i-\varpi]$ for $\varpi\le i<|w|$.
For a set $L$ of words, $\pr(L)$ is the set of prefixes of its members, and
$\pr(v)$ abbreviates $\pr(\{v\})$. For a nonempty root word $\rho$,
$\pr(\rho^*)$ thus consists of the prefixes of its finite powers; these are
exactly the finite prefixes of $\rho^\omega$.

\begin{lemma}[Periodic prefixes]\label{lem:periodic}
For $\rho\ne\eps$, $u\in\pr(\rho^*)$ iff $u\pre\rho^\omega$ iff $u\pre\rho u$.
\end{lemma}
\begin{proof}
The first equivalence follows by choosing a sufficiently long power. If
$u\pre\rho^\omega$, then both $u$ and $\rho u$ prefix the same infinite word.
Conversely, $u\pre\rho u$ either gives $u\pre\rho$, or lets us cancel an
initial $\rho$ and repeat on a shorter word.
\end{proof}

The \emph{polarity} of a path counts contravariant positions modulo two:
\begin{equation}\label{eq:g1}
 p(w)=\bigoplus_{j<|w|}c(w[j]),\qquad p(\eps)=0,\qquad p(uv)=p(u)\oplus p(v).
\end{equation}
Occurrences are counted with multiplicity: if $c(0)=1$ and $c(1)=0$, then
$p(0)=1$, $p(00)=0$, and $p(010)=0$.

A \emph{tree} is a partial labeling $t:\Dir^*\rightharpoonup\{\Bot,f,\Top\}$,
with a defined root and prefix-closed domain, such that
\begin{equation}\label{eq:g2}
 wi\in\operatorname{dom}(t)\quad\Longleftrightarrow\quad t(w)=f
 \qquad(i\in\Dir).
\end{equation}
The \emph{extreme labels} are $\bot$ and $\top$. They label leaves;
an $f$-node has $n$ children, so $f()$ is also a leaf at arity zero.
Write $f(t_i)_{i\in\Dir}$ for a constructor tree. A subtree at $\eta$ has
labeling $w\mapsto t(\eta w)$, when $\eta$ belongs to the domain.
For example, the finite binary tree $f(\Bot,\Top)$ has domain
$\{\eps,0,1\}$. The full constructor tree $F=f(F,\ldots,F)$ is infinite
when $n>0$ and is $f()$ when $n=0$.

Let $\Trees$ contain all trees, $\Trees_{\mathrm{fin}}$ those with finite
domain, and $\Trees_{\mathrm{reg}}$ those with finitely many distinct rooted
subtrees. A tree is regular iff it is the unfolding of a finite labeled graph
with $n$ ordered outgoing edges at each constructor node; for the forward
direction, take the graph of its finitely many distinct subtrees. The tree $F$ has a one-state
graph. Every finite tree is regular, but $F$ shows that the converse fails
at positive arity. At arity zero, the only trees are $\Bot,f(),\Top$.

\subsection{Non-Structural Subtyping}
Order the labels by $\Bot<f<\Top$. Define
\begin{equation}\label{eq:g3}
 t\le_c u\quad\Longleftrightarrow\quad
 \forall w\in\operatorname{dom}(t)\cap\operatorname{dom}(u),\quad
 \begin{cases}t(w)\le u(w),&p(w)=0,\\u(w)\le t(w),&p(w)=1.\end{cases}
\end{equation}
For $c=0$, write $\le$ and call the order \emph{covariant}.
The order is non-structural because $\Bot\le_c t\le_c\Top$ holds even for
trees of different shapes. For binary products $c(0)=c(1)=0$; for arrows
$c(0)=1,c(1)=0$, and the usual rule is
$f(a_0,a_1)\le_c f(b_0,b_1)$ iff $b_0\le_c a_0$ and $a_1\le_c b_1$.
For instance, $f(\Top,\Bot)\le_c f(\Bot,\Top)$ for arrows.

For a relation $R$ on trees, define the textbook one-step rule
\begin{equation}\label{eq:g4}
\begin{split}
 \operatorname{Step}_c(R)(t,u)\quad\Longleftrightarrow\quad& t=\Bot\ \lor\ u=\Top\\
 &\lor\ \exists\vec a,\vec b\;[t=f(\vec a)\land u=f(\vec b)\land
 \forall i\;\begin{cases}R(a_i,b_i),&c(i)=0,\\R(b_i,a_i),&c(i)=1.\end{cases}]
\end{split}
\end{equation}
Then $\le_c=\nu R.\operatorname{Step}_c(R)$, the \emph{greatest fixed point}:
the largest relation preserved by this rule. To prove it, compare roots
in \eqref{eq:g3} and apply \eqref{eq:g1} to the children; this shows that
\eqref{eq:g3} is a fixed point. Conversely, if $R\subseteq\operatorname{Step}_c(R)$,
induction on a common path proves \eqref{eq:g3} for every pair in $R$.
A contravariant first letter exchanges the trees and toggles polarity; a
leaf case leaves only the root comparison. The greatest fixed point is
necessary: at positive arity $F\le_c F$, whereas the least fixed point
cannot justify this comparison by finitely many unfoldings. On finite
trees, the same rule is the usual inductive order. The relation is a partial
order, as also follows from normalization and the pointwise characterization
in \eqref{eq:g11}.

\subsection{Constraints, Solutions, and Entailment}
Let $V$ be a finite set of $k$ variables containing $x,y$. A \emph{flat constraint system} $\phi$ is a finite conjunction of literals
\begin{equation}\label{eq:g5}
 z\le f(a_i)_i,\qquad f(a_i)_i\le z,\qquad z=\Bot,\qquad z=\Top,
\end{equation}
where each $a_i$ is a variable. An assignment $\sigma:V\to\mathcal D$, for
$\mathcal D\in\{\Trees,\Trees_{\mathrm{reg}},\Trees_{\mathrm{fin}}\}$, is a
\emph{solution} if it satisfies all literals, interpreting inequalities by
$\le_c$. Write $\Sat_{\mathcal D,c}(\phi)$ for existence of a solution. Define
\begin{equation}\label{eq:g6}
 \phi\models_{\mathcal D,c}x\le y\quad\Longleftrightarrow\quad
 \forall\sigma:V\to\mathcal D,\quad\sigma\models_c\phi\Longrightarrow
 \sigma(x)\le_c\sigma(y).
\end{equation}
Omit $\mathcal D$ for arbitrary trees, and use $\mathrm{reg}$ and $\mathrm{fin}$
for the other domains; omit $c$ as well in the covariant case.
Satisfiability asks for one solution; entailment asks
about every solution. A \emph{countermodel} is a solution that violates
the query. An unsatisfiable system entails every query.
For example, $x\le f(a_i)_i$ and $f(a_i)_i\le y$ entail $x\le y$.
With distinct variables they do not entail $y\le x$: set $x=\Bot,y=\Top$.
For a binary covariant constructor, $f(x,x)\le x\le f(x,x)$ has a regular
infinite solution but no finite solution.
The query $x\le y$ is a variable inequality. It is accepted directly and is
not a literal of \eqref{eq:g5}; general systems and queries reduce to this
form as follows.

\begin{proposition}[Flattening]\label{prop:flattening}
At positive arity, any finite system of inequalities between finite type
expressions, with a type-expression query, reduces to the flat syntax
\eqref{eq:g5} and a variable query. Solutions correspond under projection
in each of the three tree domains. At arity zero, such entailments are
decidable by enumerating assignments in $\{\bot,f(),\top\}$.
\end{proposition}
\begin{proof}
Give each nonvariable subexpression a fresh name. For a constructor
subexpression, impose both flat bounds between its name and the constructor
applied to the names of its children; for a constant, impose its constant
literal. Antisymmetry forces each name to denote its subexpression.
Replace each original inequality, and the query, by the variable inequality
between the names of its two sides. The query is now a variable query.

A variable inequality $u\le v$ in the system is not a literal of
\eqref{eq:g5}, but at positive arity it can be expressed with fresh
auxiliaries. Choose a child $i$, put $u,v$ in the $i$th coordinates of
vectors $\vec a,\vec b$ (exchanged if $c(i)=1$), use a common filler
variable in the other coordinates, and add the two literals
$f(\vec a)\le z\le f(\vec b)$ with a fresh \emph{sandwich variable} $z$.
A solution extends by evaluating subexpressions and taking $z$ to be
$f(\vec a)$. Conversely, the two literals give $f(\vec a)\le_c f(\vec b)$,
and decomposing it at child $i$ gives $u\le_c v$.
These extensions use only constants and finitely many constructor
applications, which preserve finite and regular trees. At arity zero, all
expressions evaluate in the displayed three-element chain, and the sandwich
encoding is unavailable.\footnote{Every flat literal
then concerns one variable. Solution sets, even after existential projection,
are products of unary sets; the order on $\bot<f()<\top$ is not such a product.}
\end{proof}
This reduction explains the mathematical scope; the exported formal
procedures take the flat syntax directly (\cref{app:lean}).

\subsection{Normalization}\label{sec:normal}
Normalization moves all order reversals from the comparison rule into the
leaf labels: after it, every path is compared covariantly.
Define \emph{tree duality} $\mathfrak d(t)$ by exchanging $\bot$ and
$\top$ and leaving every $f$-label unchanged, including the nullary
constructor $f()$. For $q\in\{0,1\}$, define $H_qt$ on the same domain as
$t$, exchanging $\bot$ and $\top$ in its label at $w$ precisely when
$q\oplus p(w)=1$. The following lemma shows that $H_0$ turns $\le_c$ into
the covariant order.

\begin{lemma}[Normalization]\label{lem:normalization}
\begin{equation}\label{eq:g9}
H_q f(t_i)_i=f(H_{q\oplus c(i)}t_i)_i,\qquad
 H_1=\mathfrak d H_0,\qquad H_0H_0=\mathrm{id},
\end{equation}

\begin{equation}\label{eq:g10}
t\le_c u\quad\Longleftrightarrow\quad H_0t\le H_0u.
\end{equation}

\end{lemma}
\begin{proof} Exchanging the extreme labels preserves \eqref{eq:g2}. At a child path $iw$, \eqref{eq:g1} gives
the sign in \eqref{eq:g9}; flipping twice proves the involution identities.
At every common path, exchanging the extreme labels reverses their order. Thus the comparison in
\eqref{eq:g3} is exactly the covariant comparison after applying $H_0$, proving \eqref{eq:g10}.
In particular, $H_0$ is an order isomorphism to the covariant tree order.
Both $H_q$ and $\mathfrak d$ preserve finite domains. They preserve regularity:
in a finite graph, attach one polarity bit to each node and update it on an
edge labeled $i$ by xor with $c(i)$. \end{proof}

\subsection{Unsafe Words and Sides}\label{sec:unsafe}
For a covariant tree $t$, define its \emph{constant leaf extension}
$\widehat t:\Dir^*\to\{\bot,f,\top\}$ by following a word until its
end or the first leaf, and then retaining that leaf's label. Thus
$\widehat t(w)=\top$ exactly when some prefix of $w$ is a $\top$-leaf
of $t$; the analogous assertion holds for $\bot$.

\begin{lemma}[Leaf-extension characterization]\label{lem:leaf-extension}
The covariant order has the characterization

\begin{equation}\label{eq:g11}
t\le u\quad\Longleftrightarrow\quad
 \forall w\quad\widehat t(w)\le\widehat u(w).
\end{equation}

\end{lemma}
\begin{proof} The reverse implication follows by considering common domain
paths. For the forward implication, follow both trees along $w$. Until a
leaf is reached, labels compare by hypothesis. If the left tree stops at
$\bot$, or the right at $\top$, the extended comparison holds forever.
If the left stops at $\top$, the right must also be $\top$ at that common
node; if the right stops at $\bot$, the left must also be $\bot$.
Otherwise, both reach the end and the original comparison applies. \end{proof}

For a tree $t$ ordered by $\le_c$, put

\begin{equation}\label{eq:g12}
\operatorname{Pref}_c^+(t,w):\Longleftrightarrow\widehat{H_0t}(w)=\top,
 \qquad
 \operatorname{Pref}_c^-(t,w):\Longleftrightarrow\widehat{H_0t}(w)=\bot.
\end{equation}

We say that $t$ has a \emph{top prefix} along $w$ if
$\operatorname{Pref}_c^+(t,w)$, and a \emph{bottom prefix} along $w$ if
$\operatorname{Pref}_c^-(t,w)$.
Thus a top prefix is an actual top leaf at an even-polarity prefix of $w$,
or an actual bottom leaf at an odd-polarity prefix. Its polarity is measured
at that prefix, not at a later point past the leaf.
A word $w$ is \emph{left unsafe} for an assignment $\sigma$ if
$\operatorname{Pref}_c^+(\sigma x,w)$ and not $\operatorname{Pref}_c^+(\sigma y,w)$, and
\emph{right unsafe} for $\sigma$ if
$\operatorname{Pref}_c^-(\sigma y,w)$ and not $\operatorname{Pref}_c^-(\sigma x,w)$.
It is left (right) unsafe for $(\phi,x,y)$ if it is left (right) unsafe for
some solution of $\phi$.

By \eqref{eq:g10}--\eqref{eq:g11},

\begin{equation}\label{eq:g13}
\phi\not\models_c x\le y
 \quad\Longleftrightarrow\quad
 \text{some word is left unsafe or right unsafe}.
\end{equation}

Indeed, the only failures of order in the three-element label chain are a
left top against a non-top, or a right bottom against a non-bottom. The same
argument shows, for each assignment $\sigma$, that
$\sigma(x)\not\le_c\sigma(y)$ iff some word is left or right unsafe for
$\sigma$.


\begin{example}[Running arrow system]\label{ex:running-arrow}
Take $c=(1,0)$, variables $x,y,a,b$, and the three flat literals
\[
 \phi=\{x\le f(a,b),\ f(a,b)\le x,\ y\le f(a,b)\}.
\]
The first two force $x=f(a,b)$, so every solution has $y\le_c x$.
The converse query fails: the finite assignment
\[
 \sigma(a)=\Bot,\quad \sigma(b)=\Top,\qquad
 \sigma(x)=f(\Bot,\Top),\quad \sigma(y)=f(\Top,\Bot)
\]
solves $\phi$, but $\sigma(x)\not\le_c\sigma(y)$.
As \cref{fig:arrow-witness} shows, normalization gives
$H_0\sigma(x)=f(\Top,\Top)$ and $H_0\sigma(y)=f(\Bot,\Bot)$.
Thus $w=0$ witnesses both sides of unsafety: $x$ has a top prefix and
$y$ does not, while $y$ has a bottom prefix and $x$ does not.
The same assignment witnesses both sides for $00$: the leaves are reached
at $0$, whose polarity is $1$, even though $p(00)=0$.
We reuse this system with the failing query $x\le y$ below.
\end{example}

\begin{figure}[t]
\centering
\begin{tikzpicture}[x=10mm,y=8mm,>=Stealth,
    every node/.style={font=\small,inner sep=2pt}]
  \node (x) at (0,0) {$f$};
  \node (x0) at (-1,-1.2) {$\Bot$};
  \node (x1) at (1,-1.2) {$\Top$};
  \node[above=2mm of x] {$\sigma(x)$};
  \draw (x) -- node[left,font=\scriptsize] {$0\;(-)$} (x0);
  \draw (x) -- node[right,font=\scriptsize] {$1\;(+)$} (x1);
  \node at (2.6,-0.5) {$\not\le_c$};
  \node (y) at (5.2,0) {$f$};
  \node (y0) at (4.2,-1.2) {$\Top$};
  \node (y1) at (6.2,-1.2) {$\Bot$};
  \node[above=2mm of y] {$\sigma(y)$};
  \draw (y) -- node[left,font=\scriptsize] {$0\;(-)$} (y0);
  \draw (y) -- node[right,font=\scriptsize] {$1\;(+)$} (y1);
  \draw[->] (0,-1.8) -- (0,-2.5) node[midway,right] {$H_0$};
  \draw[->] (5.2,-1.8) -- (5.2,-2.5) node[midway,right] {$H_0$};
  \node at (0,-2.9) {$f(\Top,\Top)$};
  \node at (2.6,-2.9) {$\not\le$};
  \node at (5.2,-2.9) {$f(\Bot,\Bot)$};
\end{tikzpicture}
\caption{The countermodel of \cref{ex:running-arrow}. The argument edge $0$ (the arrow's domain)
reverses the order: $x\le_c y$ would require $\Top\le_c\Bot$ there.
Normalization flips that edge's leaf labels, exposing the top-prefix and
bottom-prefix failures at the same word $0$.}
\Description{Two binary constructor trees. The x tree has bottom on child zero
and top on child one; the y tree has top on child zero and bottom on child one.
Child zero is contravariant and child one covariant. Arrows below show
normalization to two top leaves for x and two bottom leaves for y.}
\label{fig:arrow-witness}
\end{figure}


\FloatBarrier
\subsection{Notation at a Glance}\label{sec:notation}
The following table collects the main notation, including symbols
introduced later, grouped by the part of the paper that introduces them.
In particular, the three dual operations have different input and output
types.
\begingroup
\small
\begin{longtable}{@{}>{\raggedright\arraybackslash}p{.21\linewidth}>{\raggedright\arraybackslash}p{.28\linewidth}>{\raggedright\arraybackslash}p{.44\linewidth}@{}}
\toprule
Symbol & Type & Reading or effect\\
\midrule
\endfirsthead
\toprule
Symbol & Type & Reading or effect\\
\midrule
\endhead
\bottomrule
\endfoot
\multicolumn{3}{@{}l}{\emph{Semantics and constraints (\cref{sec:background,sec:variance,sec:spines})}}\\
$\Trees$, $\Trees_{\mathrm{reg}}$, $\Trees_{\mathrm{fin}}$ & sets of trees & Arbitrary, regular, and finite trees; $\mathcal D$ ranges over these domains.\\
$V$, $V^\pm$, $\sigma$ & variable sets; assignment & Original and signed variables; a tree for each variable.\\
$\phi$, $\psi$, $\chi$ & finite constraint systems & Original system, signed translation, and a generic covariant system.\\
$k$ & integer & Number of input variables (never a block count).\\
$\Dir$, $c$, $p$ & alphabet; $\Dir\to\{0,1\}$; $\Dir^*\to\{0,1\}$ & Child positions, their variances, and path polarity.\\
$w[i]$, $w[:j]$, $w[j:]$ & letter; word; word & Letter at position $i$; prefix and suffix at cut $j$.\\
$u\pre w$, $\pr(L)$ & relation on words; set of words & Prefix; all prefixes of words in $L$, with $\pr(v)=\pr(\{v\})$.\\
$\mathfrak d$ & $\Trees\to\Trees$ & Tree duality: exchange $\bot,\top$, retain $f$.\\
$H_q$ & trees under $\le_c$ $\to$ covariant trees & Normalize labels using $q\oplus p(w)$.\\
$\operatorname{Pref}^\pm_c(t,w)$ & truth value & Top ($+$) or bottom ($-$) prefix of $t$ along $w$.\\
$\flip$ & $V^\pm\to V^\pm$ & Variable flip: complement the sign.\\
$S$ & signed assignments $\to$ signed assignments & Assignment involution: $(S\sigma)(v)=\mathfrak d(\sigma(\flip v))$; fixed witnesses are $S$-fixed.\\
$\theta$; $X_\theta$, $Y_\theta$ & side; signed variables & Left or right side and its two query roots in $\psi$: $(x^+,y^+)$ on the left, $(y^-,x^-)$ on the right.\\
$D_\chi$ & relation on variables & Reflexive, transitive decomposition closure.\\
$\Up$, $\Lo$ & word $\times$ variable $\times$ variable $\to$ truth value & Supported upper and lower path bounds.\\
$\emb_b(\eta,t)$ & tree & Embed $t$ at path $\eta$, with off-path filler $b$.\\
$\ell_f,u_f,u_\bot$ & predicates on pairs of a variable and a path & Lower constructor, upper constructor, and upper bottom support.\\
$\Psi^4_w$; $P_j$, $Q_j$, $Y_j$, $W_j$ & system; signed variables & Four-spine extension along $w$ and its spine variables at cut $j$.\\
$U_j,L_j$; $U^R_j,L^Z_j$ & sets of signed variables & Variables reached by upper and lower path bounds at cut $j$, from $x^+$ and $y^+$, or from the spine roots $R_0$ and $Z_0$.\\
\midrule
\multicolumn{3}{@{}l}{\emph{Readers and holes (\cref{sec:holes,sec:tail,sec:depth,sec:decision})}}\\
$(\mathcal M,\Adm,\Rej)$, $N$ & reader; integer & Deterministic automaton with admission relation and rejected states; its number of states.\\
$\mathcal M_\theta$, $\Adm_\theta$, $\Rej_\theta$ & reader & Left or right label reader: the label reader of \cref{sec:label-reader} on side $\theta$.\\
$E_i$, $G_i$ & maps on sets of signed variables & One-letter updates of the upper and lower sets.\\
$(s,e)$ & pair of cuts & Comparison when $s<e$ and $w[e:]\pre w[s:]$.\\
$q_j$, $\Rej$, $\Adm$ & state; set of states; relation on states & Reader run, rejection, and admission.\\
$J(w)$, $r(w)$ & integers & First rejected cut and rejected-tail length $|w|-J(w)$.\\
hole; fixed witness & word; assignment & Rejected word without admitted comparison, for a satisfiable system an unsafe word; $S$-fixed solution showing unsafety along a word. \emph{Witness} always means the latter.\\
$B$, $d$, $H$ & integers & Tail bound $4N-2$, depth bound $2N(B+1)-1$, and length bound $\holeBound(N,d)$ (\cref{thm:hole-decision}).\\
\midrule
\multicolumn{3}{@{}l}{\emph{Hierarchies and compression (\cref{sec:depth,sec:search,sec:complexity})}}\\
$u_j$, $\ell_j$, $\Cores_j$, $\Supp_j$ & word; letter; set of cuts; set of states & Hierarchy level, marker, ordinary cores, and their reader states (state support).\\
$\mathcal M_j$, $\triangleright$, $\star$ & lettered monitor; label; state & Derived monitor, start label, and new derivative start state.\\
$\Tfin$, $\Lambda$, $\Lfin$ & predicate; two relations on labels & Final target, nonterminal admission, and terminal admission.\\
$\mathcal A_i$, $g_i$ & set of letters; integer & Level-$i$ dictionary (letters of $u_i$) and its size.\\
$\Pi_i$, $Z_i$ & set of letters; set of states & Protected letters and monitor states in use at level $i$.\\
$\mathcal E_j$ & set of letters & Letter support kept at level $j$ during shortening.\\
$\Delta$ & integer & Depth of the hole being shortened in \cref{sec:search}.\\
$\TyBound(N,d)$, $\AlBound(N,d)$, $\DictBound(N,d)$ & integers & Number of exchange types, allowance, and dictionary bound \eqref{eq:dictionary-constants}.\\
$h$, $\mathit{col}(N,h)$ & integers & Number of blocks (tracks) shortened together; column count of their summaries \eqref{eq:col}.\\
$\holeBound(N,d)$ & integer & Length bound for holes of depth at most $d$ \eqref{eq:holebound}.\\
dictionary-normal & property of holes & Normal form of \cref{lem:dictionary-normal-form}; unrelated to the normalization $H_q$.\\
\midrule
\multicolumn{3}{@{}l}{\emph{Complexity (\cref{sec:complexity})}}\\
$\insize{\phi}$ & integer & Size of the explicit input.\\
cycle clash & property of a system & $\Lo(\eta,a,a)$, $D_\chi(a,b)$ and $\Up(\eta,b,b)$ for a nonempty $\eta$; excludes finite solutions (\cref{sec:decision}).\\
$\mathcal W$, $\Gamma_{\mathcal W}$, $\Theta$ & lettered monitor; alphabet; summary & Current monitor of the summary search, its abstract letters $1,\ldots,h$, and a column summary determining the next monitor.\\
$\mathbb T$ & alphabet & Counter tokens $\{0,1,\#,\$\}$ of the long-hole construction.\\
$T_P$, $T_P^m$ & trees & Top-free tree with constructor positions $P$, and its truncation at depth $m$ with $\Top$ at the cut.\\
$\Phi^{(n)}(\mathcal A)$, $\mathcal A^\circ$ & system; automaton & Rejection system of an all-final automaton $\mathcal A$; all-final code of an automaton (hardness).\\
token encoding; $e$ & maps into $\{0,1\}^*$ & Binary encoding of counter tokens (long holes); letter doubling $a\mapsto0a$ (hardness).\\
\end{longtable}
\endgroup
